
\section{Conclusion}

This paper introduces the background technology of our proposed deep predictive learning, such as predictive coding, related research, and robot applications. We also describe our approach and application examples for complex robot motion generation using deep predictive learning from the following three perspectives.


\begin{enumerate}
\renewcommand{\labelenumi}{\roman{enumi}}
\item Switching multiple prediction models based on prediction error.
\item Combining multiple motions by designing the dynamics inside the model.
\item Acquiring integrated representation of language-motion and its application to complex motion generation.
\end{enumerate}


Although each method has been developed separately, they share many common elements and are expected to be integrated in the future. However, some issues need to be addressed, such as training data collection and segmentation methods, to improve accuracy and apply the methods to challenging tasks. In addition, deep predictive learning highly depends on the learned physical characteristics. This can be addressed by transfer learning or other methods when there is some similarity in physical characteristics, but it cannot be applied to robots with completely different structures. We plan to revisit the concept of deep predictive learning from the perspective of cognitive robotics, where the concept originated, and apply the method to a longer-term and more diverse set of tasks better suited to real-world problems.
